\subsection{Regular sequences and scalar extension}

PROPOSITION 5.--- Let \(\mathfrak{p}: A \to B\) a local homomorphism of Noetherian local rings, N a finitely generated B-module, \(\mathbf{x} = (x_1, \ldots, x_r)\) be a sequence of elements of \(\mathfrak{m}_B\) and \(u: A[T_1, \ldots, T_r] \longrightarrow B\) the unique homomorphism of A-algebras such that \(u(T_i) = x_i\) for \(i = 1, \ldots, r\) The following conditions are equivalent:

\begin{enumerate}
\def\labelenumi{(\roman{enumi})}
\item
  the homomorphism u makes N an A{[}T \(_{1}\) ,\ldots{} ,T \(_{r}\) {]}-module flat ;
\item
  the A-module N is flat and, for every A-module M, the sequence x is M ⊗A N-regular;
\item
  the A-module N is flat and the sequence x is \(\kappa_{A} \otimes_{A} N\)-regular ;
\item
  the A-module N/(x₁N + \ldots{} + xᵣN) is flat and the sequence x is N-regular.
\end{enumerate}

(i)⇒(ii): set \(\mathbf{T} = (\mathrm{T}_{1}, \ldots, \mathrm{T}_{r})\) Suppose that the A{[}T{]}-module N is flat. Since A{[}T{]} is flat over A, the A-module N is flat (I, § 2, no. 7, cor. 3 of prop. 8). Let M be an A-module. The sequence T is evidently M{[}T{]}-regular, hence M{[}T{]} ⊗\_\{A{[}T{]}\} N-regular, since N is flat over A{[}T{]}Now the A{[}T{]}-module M{[}T{]} ⊗\_\{A{[}T{]}\} N is identified with M ⊗\_\{A\} N, the homothety of ratio T\_i corresponding to the endomorphism 1\_M ⊗ (x\_i)\_N. Condition (ii) is therefore satisfied.

\begin{enumerate}
\def\labelenumi{(\roman{enumi})}
\setcounter{enumi}{1}
\tightlist
\item
  \(\Rightarrow\) (iii): this is trivial.
\end{enumerate}

(iii)⇒(iv): this follows from prop. 10 of § 1, no. 6.

(iv)⇒(i): let t denote the ideal of A{[}T{]} generated by T. The (A{[}T{]}/t)-module N/tN is flat by hypothesis, and N is ideally separated for t (III, § 5, no. 4, prop. 2). To prove that N is flat over A{[}T{]}, it suffices to prove, in view of th. 1 of loc. cit., no. 2, that the A-module \(\operatorname{Tor}_{1}^{A{[}T{]}}(A{[}T{]}/t,N)\) is zero. But since the sequence T is A{[}T{]}-regular, this module is isomorphic to \(\mathrm{H}_{1}(T,N)\) (A, X, p.~159, remark 3), which is zero since the sequence T is N-regular (loc. cit., p.~157, prop. 5).

COROLLARY. Let \(k\) a field, \(\Lambda\) an \(k\) -local Noetherian algebra, \(\mathbf{x} = (x_1, \ldots, x_r)\) be a sequence of elements of \(\mathfrak{m}_{\mathrm{A}}\) and M a \(\Lambda\) finite type module. Let \(\widehat{\mathrm{A}}\) and \(\widehat{\mathrm{M}}\) the completions of A and M for their topology ( \(x_1\mathrm{A} + \ldots + x_r\mathrm{A}\) -adic; denote by \(u: k[\mathrm{T}_1, \ldots, \mathrm{T}_r] \longrightarrow \mathrm{A}\) the unique homomorphism of \(k\) -algebras such that \(u(\mathrm{T}_i) = x_i\) for \(i = 1, \ldots, r\) , and by \(\hat{u}: k[[\mathrm{T}_1, \ldots, \mathrm{T}_r]] \longrightarrow \widehat{\mathrm{A}}\) the unique continuous homomorphism extending it. The following conditions are equivalent:

\begin{enumerate}
\def\labelenumi{(\roman{enumi})}
\item
  the sequence x is M-regular;
\item
  the homomorphism u makes M into a \(k[T_{1},\ldots,T_{r}]\) -module flat ;
\item
  the homomorphism \(\hat{u}\) makes \(\hat{M}\) a \(k[[T_{1},\ldots,T_{r}]]\) -module flat.
\end{enumerate}

The equivalence of (i) and (ii) follows from the equivalence of conditions (i) and (iv) of Prop. 5; the equivalence of (ii) and (iii) follows from III, § 5, no. 4, Prop. 4.

These results make it possible to characterize Macaulay modules in two important cases. Let us denote \(\Lambda\) a Noetherian local ring, M a finitely generated A-module. It is equivalent to say that the \(\Lambda\) -module M is Macaulay or that the \(\widehat{A}\) -module \(\widehat{M}\) is Macaulay \(\S\) 2, no. 7, cor. 4 of prop. 8). We shall henceforth assume that the Noetherian local ring A is complete.

\begin{enumerate}
\def\labelenumi{\arabic{enumi})}
\tightlist
\item
  Suppose first that A has a subfield; it then admits a field of representatives. \(k\) (IX, § 3, no. 3, th. 1). Let \((x_{1},\ldots ,x_{r})\) a maximal secant sequence for M; denote \(u:k[[\mathrm{T}_{1},\dots ,\mathrm{T}_{r}]]\longrightarrow \Lambda\) the unique continuous homomorphism such that \(u(\mathrm{T}_i) = x_i\) for \(i = 1,\dots ,r\) By lemma 4 b) of IX, § 2, no. 5 and remark 1 of VIII, § 3, no. 2, A/Ann(M) is a \(k[[\mathrm{T}_{1},\dots ,\mathrm{T}_{r}]]\) -module of finite type and consequently M is a \(k[[\mathrm{T}_{1},\dots ,\mathrm{T}_{r}]]\) -finite type module. That said, the following conditions are equivalent:
  \begin{enumerate}
  \def\labelenumii{(\roman{enumii})}
  \item
  the \(k[[\mathrm{T}_1,\dots ,\mathrm{T}_r]]\) -module M is free;
\item
  the A-module M is Macaulay.
\end{enumerate}

Indeed, it is equivalent to say that M is a Macaulay A-module or that the sequence \((x_{1},\ldots,x_{r})\) is M-regular (§ 2, no. 3, th. 1). By the above corollary, this latter condition means that \(k[[T_{1},\ldots,T_{r}]]\) the module M is flat, or equivalently that it is free since it is finitely generated.

\item
  Suppose that the residue field \(\kappa_{\Lambda}\) of A be of characteristic \(p > 0\) and that we have \(\dim(M/pM) < \dim(M)\) . Let \((x_1, \ldots, x_r)\) a maximal secant sequence for M/pM, so that \((p1_{\mathrm{A}}, x_1, \ldots, x_r)\) is a maximal secant sequence for M. Let C be a \(p\) -ring of length \(+\infty\) with residue field \(\kappa_{\Lambda}\) (IX, § 2, no. 3, prop. 5). There exists a homomorphism \(u_0\) from C into A which induces the identity on the residue fields; let \(u: \mathrm{C}[[\mathrm{T}_1, \ldots, \mathrm{T}_r]] \longrightarrow \mathrm{A}\) the unique homomorphism extending \(u_0\) and mapping \(\mathrm{T}_i\) on \(x_i\) for every \(i\) It follows as before from loc. cit., no. 5, lemma 4 that \(u\) makes M a \(\mathrm{C}[[\mathrm{T}_1, \ldots, \mathrm{T}_r]]\) -module of finite type. The following conditions are equivalent:
  \begin{enumerate}
  \def\labelenumii{(\roman{enumii})}
\item
  the C{[}{[}T1, \ldots, Tr{]}{]}-module M is free;
\item
  the A-module M is Macaulay.
\end{enumerate}

Indeed, condition (ii) is equivalent to saying that the sequence \((x_{1},\ldots,x_{r})\) is M-regular and that the homothety with ratio p in \(\mathrm{M}/(x_{1}\mathrm{M}+\ldots+x_{r}\mathrm{M})\) is injective (§ 2, no. 3, th. 1). But this latter condition means that the C-module \(\mathrm{M}/(x_{1}\mathrm{M}+\ldots+x_{r}\mathrm{M})\) is torsion-free, hence flat (Λ, X, p.~9, example 7). Thus, in view of prop. 5, (iv)⇔(i), condition (ii) is equivalent to the fact that the \(C[T_{1},\ldots,T_{r}]\) that the module M is flat, or equivalently (III, § 5, no. 4, prop. 4) that the \(C[[T_{1},\ldots,T_{r}]]\) -module M is flat, that is, free since it is finitely generated.

\end{enumerate}

